A colored image is composed of three channels: red, green, and blue, each of which has a value range from 0 to 255. By subtracting one value of a pixel by 1, the image can be considered as changed and be read by a model as new data. That means there are unlimited image augmentation techniques. The challenge now is how to select efficient augmentation methods. This question is  an active topic in Computer Vision as the field of machine learning is continuously exploring new ideas and challenging existing techniques. 

In this section, we aim to provide an overview of the current state-of-the-art augmentation techniques. To enhance readability and comprehensibility, we have grouped these techniques into two main categories: data warping, oversampling and undersampling.

\subsection{Data warping}
One of the earliest and most commonly used image augmentation techniques is data warping. Data warping simply means transforming existing images to create new ones. This includes augmentations such as geometric transformations (e.g. flipping, rotation, and translation), color transformations (e.g. blocking one or two of the red, green, and blue channels), and random crop (cropping patches from each image). These techniques are relatively simple to implement and have been shown to be effective in improving the performance of machine learning models on many data-sets \cite{shorten2019survey}. However, there are several more advanced data warping techniques used in augmentation beyond basic ones such as AutoAugment, RandAugment, Cutout, Random Erasing, Hide and Seek and Grid Mask.

\subsubsection{AutoAugment, Fast AutoAugment, RandAugment}
AutoAugment is a relatively recent image augmentation technique that uses reinforcement learning to find the optimal set of augmentations to apply to an image \cite{cubuk2018autoaugment}. This augmentation technique tackle the questions "if there are many augmentations, which ones should be implemented and with what maginitude?". The authors made a space for searching where there are different sets of rules, each with a way of changing an image, like flipping or changing colors. To find the best combination of these rules, a reinforcement learning approach is implemented to navigate this search space. The ideal augmentation policy is determined by the probabilities and magnitudes of the functions. While AutoAugment can produce effective policies, it is computationally expensive and requires significant resources, such as 1500 GPU hours to explore a good policy for the ImageNet dataset \cite{lim2019fast}.

In an effort to address this limitation, Lim et al. proposed Fast AutoAugment \cite{lim2019fast}, which reduces the search phase and allows for adjustment based on different model sizes. This results in a more efficient and less computationally intensive method with improved regularization strength, requiring less than one-third of the time required by AutoAugment. 

A year after the creation of AutoAugment, Cubuk et al. introduced RandAugment \cite{cubuk2020randaugment}, which addressed two of AutoAugment's disadvantages. The separate search phase of AutoAugment incurred increased computational cost and weak regularization strength for varying model sizes. RandAugment is a similar technique, but it gives up reinforcement learning completely. Instead, RandAugment uses random sampling. This allows RandAugment to be applied directly to the target task and be adjusted according to different model sizes. RandAugment is extremely fast and efficient. It is still being used as common practice today.

\subsubsection{Cutout, Random Erasing, Hide and Seek, GridDropout}
Recently, inspired by the effectiveness of the dropout layers in convolutional networks \cite{srivastava2014dropout}, image augmentation techniques such as Cutout, Random Erasing, and Hide And Seek are extensively researched. Cutout, Random Erasing, and Hide And Seek are image augmentation techniques that involve removing parts of images to create new ones. Cutout involves blacking out square regions of an image \cite{devries2017improved}, while random erasing covers certain parts of an image with noise, and the size and number of parts are randomly chosen \cite{zhong2020random}. Hide and seek involves hiding and then revealing parts of an image, essentially making the classifier focus on different parts of the original image \cite{singh2018hide}. A more systematic drop-out strategy is Gridmask. GridMask masks out rectangular regions of an image in a grid pattern, forcing the model to learn to recognize objects and features in different positions within the image \cite{chen2020gridmask}. Similar to weight regularization in dropout layers, these techniques encourage the model to learn from the surrounding context of an image instead of focusing solely on specific features. By randomly removing parts of an image, the model is encouraged to be more robust and better equipped to handle missing information. This is particularly useful when dealing with real-world scenarios where images may not always be complete or clear, as increasing the diversity of the training set can lead to more accurate and reliable models.

\subsubsection{Adversarial attacks}
Another approach to make the dataset more robust is to attack its weakness, such as Adversarial training. Adversarial attacks are techniques used to manipulate input data in such a way that it is misclassified by the model. So adversarial training works by introducing carefully crafted adversarial attacks into the training data of the model, which causes mislabelling by the model. The model will be punished and then learns to recognize these adversarial examples and adjust its parameters accordingly, in order to improve its performance on both clean and adversarial data \cite{shrivastava2017learning}. 

\subsubsection{Neural style transfer}
The most recent and novel approach is augmentation using neural style transfer \cite{gatys2015neural}. This is a type of deep learning-based augmentation that changes the look of an image by applying a certain style. To use neural style transfer for image augmentation, we first need to train a style transfer model on a set of reference images with different styles. During training, the model learns how to extract the style of an image and apply it to a different image while preserving the content. However, selecting the appropriate style for image generation can be a challenging task, often requiring a significant amount of computational resources. As a result, achieving high-quality and realistic images can be quite difficult.

\subsection{Oversampling and Undersampling}
Apart from data warping, another group of image augmentation techniques involves oversampling and undersampling. The common of this group is not to transform an existing image, but to create new images based on existing ones. While oversampling refers to the creation of new images from existing ones by adding noise or altering specific features, undersampling involves removing parts of images to create new ones. These techniques can be used to balance the number of samples in different classes of a data-set, which can be useful in preventing class imbalance problems. This group of data augmentations includes mixing images, feature space augmentation, and using AI networks e.g GAN-based augmentation.

\subsubsection{Mixing Images}
We can also combine different images in the same class to do augmentation. Mixing images is a type of data augmentation that combines two or more images to create a new image \cite{summers2019improved}. There are different methods for mixing images, including simple blending, weighted blending, and alpha blending. Simple blending involves overlaying two images on top of each other using a basic blending function. Weighted blending assigns different weights to each image based on their importance or relevance to the task at hand. Alpha blending is a more sophisticated technique that involves creating a transparent mask that blends two images based on their pixel intensities.

This new image is then used as additional training data. This helps to increase the diversity of the data and can lead to better performance of the model being trained. By mixing images after image processing, the error rate on the CIFAR-10 data-set was reduced from 5.4\% to 3.8\% \cite{liang2018understanding}. However, the problem with Mixing images is its randomness. It combines pixel values from different images without taking into account the important features, which can result in the inclusion of unwanted noise from the original images. 

\subsubsection{Augmentations in latent space}
All the discussed data augmentation techniques mentioned above are used on the original images before they are processed. But by using neural networks, we can reduce the size of images and augment directly the representation in latent-space. Autoencoders have two parts, an encoder that simplifies the image, and a decoder that rebuilds the image back to its original size. By simplifying the image, the autoencoder can find the important parts and relationships between the pixels, which can then be used to make changes to the image and improve its accuracy. This information can then be used to enhance the features in the latent space, making it a useful tool for data augmentation. For examples, DeVries and Taylor used k-nearest neighbors to generate new images in the latent space and saw positive results on five different types of data sets \cite{terrance2017dataset}.

\subsubsection{GAN-based Augmentations }
Latent space augmentation can sometimes result in features with too much noise. Generative Adversarial Networks (GAN), introduced by Ian Goodfellow in 2014 \cite{goodfellow2020generative}, is another exciting strategy for data augmentation. GAN-generated methods can help in augmentation by generating synthetic data that can be used to supplement the original training data, especially in cases where the training data is limited or imbalanced. The GAN model can learn the underlying distribution of the training data and generate new data samples that are similar to the real data but have different features that can help the model generalize better.

GAN-based data augmentation involves two parts: a generator and a discriminator. The generator is trained to create images from input that the discriminator can't identify. The discriminator then tries to recognize the labels of the generated images. Eventually, the generator gets better and creates images that keep the original label even after being augmented. A vanilla GAN is refered to the GAN architecture which uses multilayer perceptron networks in the generator and discriminator networks. GAN-based data augmentation is an ongoing area of research, with various modifications made to the GAN framework such as DCGANs \cite{radford2015unsupervised}, which utilize convolutional neural networks as the architecture, Wasserstein GAN (WGAN) \cite{arjovsky2017wasserstein} which uses Wasserstein distance as loss function or CycleGANs \cite{zhu2017unpaired}, which use two GANs working in tandem to match two domains to each other without paired training data.

Chou's work in 2019 demonstrated that using a simple GAN structure for data augmentation improved the performance of sorting the defective beans task, outperforming HCADIS and YOLOv3 \cite{chou2019deep}. Similarly, in 2021, Dimoiu et al. used conditional GANs to enhance the performance of aerial image segmentation in a real context of a rural zone in Romania, with their proposed scheme outperforming other GAN implementations \cite{dimoiu2021improved}. In 2019, Zhang and colleagues used a Wasserstein GAN to handle unequal image data in detecting weld defects \cite{zhang2019weld}. Both of their experiments on augmented images and real world defect images achieve satisfying accuracy, which substantiates the possibility that the proposed approach is promising for weld defect detection. 

GANs have the potential to be very powerful, but GANs are computationally expensive and require a significant amount of resources, such as high-end GPUs and large amounts of data. This makes it difficult for researchers with limited resources to train and generate high-quality images. Addtionally, GANs may not always be effective as an augmentation tool because GAN-generated data may not cover the entire distribution of the original data. This can result in overfitting to the limited variation present in the GAN-generated images, which can lead to poor generalization performance, particularly when there are differences in the distribution between the training and test sets. These factors can make it challenging to use GANs effectively for data augmentation purposes.

\subsubsection{Downsampling}
In contrast to generating new data, downsampling presents an alternative approach for addressing the issue of an unbalanced dataset. Downsampling involves randomly deleting samples from the majority class with sufficient samples, in order to achieve an equal number of samples across all classes, or reducing the image resolution to expedite processing time. In the former case, this approach helps balance the number of samples across the classes, thereby mitigating biases towards dominant classes. In the latter, downsampling can prove useful in scenarios where the original image resolution is excessively high, leading to computationally intensive processing times. Reducing the image size enables significant processing time reduction, while retaining the essential image features. Although this method is easy to implement and helps with the computational feasibility, there is a high risk that the deleted samples contain key informative data.

Despite significant progress in the field of image augmentations, there remain several limitations that must be addressed. One such limitation is the need for a dataset-specific augmentation policy. It can be challenging to develop a strategy that is effective for all types of images, given the inherent variations and invariances that may exist within each dataset. For instance, in the case of aerial image datasets, issues such as center biases and non-uniform illuminations must be carefully considered and understood to develop effective augmentation strategies.

    
        